\chapter{Discussion}
\label{ch:discussion}

\section{The most useful result was the one that failed to appear}

The fusion measurement is the part of this project I would keep if I could keep
only one thing. The expectation of ten to twenty percent was not arbitrary; it
is roughly right for a memory bound kernel, and removing a $24N$ byte round trip
per step sounds substantial when stated in bytes. It is worth
\FusionGainPercentAtMaxN{} percent.

What the expectation missed is arithmetic intensity. The all pairs sum performs
$O(N^2)$ arithmetic on $O(N)$ data, so the ratio of operations to bytes grows
without limit as the problem grows. At \LadderMaxN{} particles the force
evaluation performs about $4 \times 10^{11}$ floating point operations while
touching a working set of a few megabytes. Removing memory traffic from a
kernel in that regime is removing something that was not costing anything: the
counters put device memory throughput at \TiledDramPct{} percent of peak.

The general lesson is that an optimisation should be sized against the resource
it relieves before it is implemented, not after. Two minutes of arithmetic with
the measured bandwidth figure would have predicted
\FusionCeilingPercent{} percent as an upper bound and saved the effort of
being surprised. That said, the fused kernel is not wasted: it is the correct
choice at small particle counts, where it wins \FusionGainPercentAtMinN{}
percent because launch overhead is then a real fraction of the step, and it
halves the memory footprint of a step.

\section{Tiling is not a bandwidth optimisation here}

The second correction is subtler. Shared memory tiling is conventionally
introduced as a way of reducing device memory traffic, and the introduction to
this report frames it that way. On this problem the framing is wrong, and the
counters say so: device memory throughput is under one percent of peak both
before and after the change. Tiling still delivers a factor of
\TilingSpeedupAtMaxN{}, but the resource it relieves is the first level cache
and the load store pipeline, not the memory bus. Global load sectors fall by a
factor of \SectorReduction{}, and the naive kernel's first level cache
throughput of \NaiveLoneTexPct{} percent of peak identifies it as the true
bottleneck.

This matters practically. Had I tried to improve the naive kernel by attacking
device memory traffic, for instance by compressing the state or by adding
prefetching, nothing would have happened. The profile pointed at the right
resource; reasoning from the textbook description would not have.

\section{Where the specification and the measurement disagreed}

Three points in the original design were changed or contradicted by
measurement, and each is recorded rather than quietly absorbed.

The host compiler floor of GCC 16.1 is not achievable with CUDA 13.3, which
refuses any GCC newer than 15. GCC 15.2 supplies the C++23 host standard the
floor was really for.

The tile width of thirty two costs \TileWidthGainPercent{} percent against a
width of \BestTileWidth{} on this card. The default was kept so the ladder
would measure the design as specified, and the sweep is reported so the choice
is informed.

The estimated runtimes in the specification are far too long. The full
validation suite, which was estimated at ten to twenty minutes on the grounds
that the ten thousand step orbit would dominate, completes in about five
seconds. The estimate assumed the orbit run would be expensive; with two
particles it is trivially cheap, and the cost is dominated instead by the ten
thousand particle Plummer runs. The full benchmark sweep, estimated at fifteen
to thirty minutes, completes in under two minutes.

\section{Limitations}

The all pairs solver is $O(N^2)$ and always will be. At two hundred thousand
particles a step costs \TiledMsAtMaxN{} milliseconds; at two million it would
cost a hundred times that. The range where this approach is the right one is
bounded, and beyond it a tree code is not an optimisation but a different
algorithm.

The fluid solver is weakly compressible, so it enforces incompressibility only
through a stiff equation of state. That bounds the time step: the sound speed
implied by the stiffness constant sets a Courant condition far tighter than the
fluid velocity would. The dam break runs at a time step of $10^{-4}$ seconds
for that reason, so the interactive step rates reported in
Section~\ref{sec:sphresults} are steps per second and not seconds of simulated
time per second.

Single precision is a deliberate choice and it has a cost. The momentum gate
measures it directly: the residual is exactly the accumulation of rounding in
the force sum. For the problems here that cost is far inside the tolerances,
but a long term integration of a planetary system would need double precision
positions, and the ladder would look different because double precision
arithmetic is much slower relative to memory on consumer hardware.

The benchmark measures the fluid from a configuration that begins as a settled
block and evolves for a few hundred steps during measurement. A fully developed
splash has a different spatial distribution and therefore a different neighbour
list length, so the fluid timings should be read as representative of the early
collapse rather than of every phase of the flow.
